\documentclass[11pt,a4paper]{article}

\usepackage[margin=1.02in]{geometry}
\usepackage{amsmath,amssymb,amsthm,mathtools}
\usepackage{bm}
\usepackage[expansion=false]{microtype}
\usepackage{enumitem}
\usepackage{array,booktabs}
\usepackage{xcolor}
\usepackage[colorlinks=true,linkcolor=blue!55!black,citecolor=blue!55!black,
            urlcolor=blue!55!black]{hyperref}

%%% notation, continuing Notes I and II
\newcommand{\Gm}{\mathbb{G}}
\newcommand{\Rm}{\mathbb{R}}
\newcommand{\Jm}{\mathbb{J}}
\newcommand{\Km}{\mathbb{K}}
\newcommand{\Ig}{\mathbb{I}}
\newcommand{\ip}[2]{\langle\!\langle #1 , #2 \rangle\!\rangle}
\newcommand{\ipG}[2]{\langle\!\langle #1 , #2 \rangle\!\rangle_{\Gm}}
\newcommand{\ipJ}[2]{\langle\!\langle #1 , #2 \rangle\!\rangle_{\Jm}}
\newcommand{\dpair}[2]{\langle #1 ; #2 \rangle}
\newcommand{\Gsharp}{\Gm^{\sharp}}
\newcommand{\Gflat}{\Gm^{\flat}}
\newcommand{\Hess}{\operatorname{Hess}}
\newcommand{\grad}{\operatorname{grad}}
\newcommand{\Jac}{\mathcal{R}}
\newcommand{\Sect}{\mathcal{K}}
\newcommand{\Damp}{\mathcal{D}}
\newcommand{\Kgain}{\mathcal{K}}
\newcommand{\Gsec}[1]{\Gamma^{\infty}(#1)}
\newcommand{\Cinf}[1]{C^{\infty}(#1)}
\newcommand{\id}{\mathrm{Id}}
\newcommand{\Lie}{\mathcal{L}}
\newcommand{\norm}[1]{\lVert #1 \rVert}
\newcommand{\normG}[1]{\lVert #1 \rVert_{\Gm}}
\newcommand{\normJ}[1]{\lVert #1 \rVert_{\Jm}}
\newcommand{\Dt}{\tfrac{D}{dt}}
\newcommand{\sym}{\operatorname{sym}}
\newcommand{\so}{\mathfrak{so}(3)}
\newcommand{\se}{\mathfrak{se}(3)}
\newcommand{\g}{\mathfrak{g}}
\newcommand{\SO}{\mathrm{SO}(3)}
\newcommand{\SE}{\mathrm{SE}(3)}
\newcommand{\ad}{\operatorname{ad}}
\newcommand{\adt}{\operatorname{ad}^{\top}}
\newcommand{\tr}{\operatorname{tr}}
\newcommand{\Ome}{\Omega_{\!e}}
\newcommand{\Omr}{\Omega_{\!r}}
\newcommand{\eR}{e_{R}}
\DeclareMathOperator{\spec}{spec}
\DeclareMathOperator{\rank}{rank}
\DeclareMathOperator*{\esssup}{sup}

\theoremstyle{plain}
\newtheorem{theorem}{Theorem}[section]
\newtheorem{proposition}[theorem]{Proposition}
\newtheorem{lemma}[theorem]{Lemma}
\newtheorem{corollary}[theorem]{Corollary}
\theoremstyle{definition}
\newtheorem{definition}[theorem]{Definition}
\newtheorem{assumption}[theorem]{Assumption}
\theoremstyle{remark}
\newtheorem{remark}[theorem]{Remark}
\newtheorem{example}[theorem]{Example}

\title{\bfseries Contraction-Based Control on $\mathrm{SO}(3)$ and $\mathrm{SE}(3)$\\[3pt]
\large Attitude control of a spacecraft, and trajectory tracking for a quadrotor}
\date{}
\author{}

\begin{document}
\maketitle
\vspace{-2.4\baselineskip}

\begin{abstract}
\noindent
Two worked designs in the framework of Notes~I and~II. For each system we
(a)~derive the dynamics on $T^{*}Q$ in canonical Hamiltonian form, reduce by
left translation to Lie--Poisson / Euler--Poincar\'e form, and identify the
affine connection control system of \cite[\S4.6.3]{BL04}; (b)~define an explicit
contraction metric and controller; (c)~prove contraction.

The spacecraft on $\SO$ is fully actuated and admits a complete, unconditional
result: the closed-loop tracking error system is exactly the shaped-and-damped
system of Note~I, and Note~I's Theorem~6.4 applies with $\mu,\sigma$ computed here
in closed form from $\Jm$, $\Km$ and the sublevel set. Two by-products:
$\Hess^{\sharp}\varphi$ is isotropised by the gain choice
$\Km = k\left(\tr(\Jm)I - 2\Jm\right)$, which is positive-definite \emph{exactly
because} the principal moments of inertia satisfy the triangle inequality
(Prop.~\ref{prop:isotropise}); and left-invariant metrics on $\SO$ have negative
sectional curvature in some planes whenever the body is inertially asymmetric,
so Note~I's velocity-bounded contraction region is a real constraint, not an
artefact.

The quadrotor on $\SE$ is under-actuated ($\rank = 4 < 6$), so Note~I's
Theorem~6.4 does \emph{not} apply. We decompose into a flat translational
subsystem --- which is isotropic and therefore contracting for \emph{every}
positive damping gain --- and the attitude subsystem of Part~A, coupled in
\emph{feedback}, not cascade. The correct tool is the small-gain lemma
\cite[Lem.~3.2]{SB13}, and the resulting theorem is explicitly conditional on two
interconnection gains that must be evaluated on the operating region.
Section~\ref{sec:audit} separates what is proved from what is assumed.
\end{abstract}

\tableofcontents

%%%%%%%%%%%%%%%%%%%%%%%%%%%%%%%%%%%%%%%%%%%%%%%%%%%%%%%%%%%%%%%%%%%%%%%%%%%%%%%
\section{Conventions}
\label{sec:conv}

Notation follows \cite{BL04} and Notes~I--II. $\hat{\cdot}:\Rm^3\to\so$ is the hat
map, $\hat x y = x\times y$, with inverse $(\cdot)^{\vee}$; $\sym(A)=\tfrac12(A+A^{\top})$.

For a Lie group $G$ with algebra $\g$ and a left-invariant metric determined by
$\Ig:\g\to\g^{*}$ symmetric positive-definite, we write
$\ip{\xi}{\eta}_{\Ig} = \dpair{\Ig\xi}{\eta}$. Two adjoints are used:
$\ad^{*}_{\xi}:\g^{*}\to\g^{*}$ is the algebraic coadjoint,
$\dpair{\ad^{*}_{\xi}\mu}{\eta} = \dpair{\mu}{\ad_{\xi}\eta}$, and
$\adt_{\xi} \coloneqq \Ig^{-1}\ad^{*}_{\xi}\Ig : \g\to\g$ is the \emph{metric}
adjoint.

\begin{lemma}[Levi-Civita connection of a left-invariant metric]
\label{lem:LC}
On left-invariant vector fields,
\begin{equation}
  \nabla_{\xi}\eta \;=\; \tfrac12\Big( [\xi,\eta] \;-\; \adt_{\xi}\eta \;-\; \adt_{\eta}\xi \Big)
  \;\eqqcolon\; \Gamma(\xi,\eta),
  \qquad\text{so}\qquad
  \nabla_{\xi}\xi \;=\; -\adt_{\xi}\xi .
  \label{eq:LC}
\end{equation}
\end{lemma}

\begin{proof}
Left-invariant functions $\ip{\eta}{\zeta}$ are constant, so the Koszul formula
reduces to
\[
  2\ip{\nabla_{\xi}\eta}{\zeta}
  = \ip{[\xi,\eta]}{\zeta} - \ip{[\eta,\zeta]}{\xi} + \ip{[\zeta,\xi]}{\eta}.
\]
Now $\ip{[\eta,\zeta]}{\xi} = \ip{\zeta}{\adt_{\eta}\xi}$ and
$\ip{[\zeta,\xi]}{\eta} = -\ip{\ad_{\xi}\zeta}{\eta} = -\ip{\zeta}{\adt_{\xi}\eta}$.
\end{proof}

For a curve with left-trivialised velocity $\xi(t)$ and any vector field $W(t)$
along it, the covariant derivative is $\Dt W = \dot W + \Gamma(\xi,W)$, so
\begin{equation}
  \Dt \xi \;=\; \dot\xi \;-\; \Ig^{-1}\ad^{*}_{\xi}\big(\Ig\xi\big).
  \label{eq:covaccel}
\end{equation}

\begin{lemma}[Left-trivialised canonical form]
\label{lem:triv}
Trivialise $T^{*}G \cong G\times\g^{*}$ by $\alpha_g \mapsto (g,\mu)$,
$\mu = T^{*}_eL_g\alpha_g$, and represent tangent vectors as $(\eta,\nu)\in\g\times\g^{*}$.
Then
\begin{equation}
  \omega_0\big((\eta_1,\nu_1),(\eta_2,\nu_2)\big)\Big|_{(g,\mu)}
  \;=\; \dpair{\nu_2}{\eta_1} - \dpair{\nu_1}{\eta_2} \;+\; \dpair{\mu}{[\eta_1,\eta_2]} ,
  \label{eq:omega-triv}
\end{equation}
and $i_{X_H}\omega_0 = dH$ with an applied body force $F\in\g^{*}$ gives
\begin{equation}
  \dot g \;=\; g\,\frac{\partial H}{\partial \mu},
  \qquad
  \dot\mu \;=\; \ad^{*}_{\partial H/\partial\mu}\,\mu \;-\; \mathbf{d}_gH \;+\; F ,
  \label{eq:LiePoisson}
\end{equation}
where $\mathbf{d}_gH\in\g^{*}$ is the left-trivialised derivative in $g$.
\end{lemma}

\begin{proof}
Matching the $\nu$- and $\eta$-components of
$\omega_0((\eta_H,\nu_H),(\eta,\nu)) = \dpair{\mathbf{d}_gH}{\eta}+\dpair{\nu}{\partial H/\partial\mu}$
gives $\eta_H = \partial H/\partial\mu$ and
$\nu_H = \ad^{*}_{\eta_H}\mu - \mathbf{d}_gH$.
\end{proof}

\begin{remark}[Why the bracket term does not contradict Note~I]
\label{rem:torsion}
Note~I, Lemma~4.1 gives $\omega_0(w_1,w_2) = \dpair{\beta_2}{u_1}-\dpair{\beta_1}{u_2}$
with \emph{no} correction term. There is no conflict: that splitting is induced by
the Levi-Civita connection, which is torsion-free, whereas the left trivialisation
\eqref{eq:omega-triv} is induced by the flat left-invariant connection
$\nabla^{-}$, whose torsion is $T(\eta_1,\eta_2) = -[\eta_1,\eta_2]$. The extra
term in \eqref{eq:omega-triv} \emph{is} that torsion, paired with $\mu$. Both
splittings are used below --- \eqref{eq:omega-triv} for the reduction, the
Levi-Civita one for the contraction metric --- and they must not be conflated.
\end{remark}

By \eqref{eq:covaccel} the forced Euler--Poincar\'e equation
$\Ig\dot\xi = \ad^{*}_{\xi}\Ig\xi - \mathbf{d}_gV + F$ is equivalent to
\begin{equation}
  \Gflat\big(\nabla_{v}v\big) \;=\; -\,dV \;+\; F ,
  \label{eq:accs}
\end{equation}
which is Note~I, eq.~(1): the affine connection control system.

%%%%%%%%%%%%%%%%%%%%%%%%%%%%%%%%%%%%%%%%%%%%%%%%%%%%%%%%%%%%%%%%%%%%%%%%%%%%%%%
\part*{Part A. Spacecraft on $\SO$}
\addcontentsline{toc}{section}{Part A. Spacecraft on $\SO$}
\setcounter{section}{0}
\renewcommand{\thesection}{A.\arabic{section}}
\setcounter{theorem}{0}
\renewcommand{\thetheorem}{A.\arabic{section}.\arabic{theorem}}

\section{Dynamics}
\label{sec:so3-dyn}

\subsection{Configuration manifold and kinetic energy metric}

$Q = \SO$, $\g = \so \cong \Rm^3$ via $\hat{\cdot}$, with
$[\hat\Omega_1,\hat\Omega_2] = (\Omega_1\times\Omega_2)^{\wedge}$. The body
angular velocity of $t\mapsto R(t)$ is $\Omega = (R^{\top}\dot R)^{\vee}$.

The inertia tensor $\Jm=\Jm^{\top}\succ0$ (body frame) defines the
\emph{left-invariant} kinetic energy metric \cite[\S5.3]{BL04}
\begin{equation}
  \Gm_R(u_R,w_R) \;=\; \Omega_u\cdot \Jm\,\Omega_w,
  \qquad \Omega_u = (R^{\top}u_R)^{\vee},
\end{equation}
so $\Ig = \Jm$ under $\so\cong\Rm^3$. We write $\normJ{x}^2 = x\cdot\Jm x$.

\subsection{Hamiltonian form on $T^{*}\SO$ and Lie--Poisson reduction}

Left-trivialise $T^{*}\SO\cong\SO\times\Rm^3$, $\alpha_R\mapsto(R,\Pi)$ with
$\Pi = \Jm\Omega$ the body angular momentum. The Hamiltonian is
\begin{equation}
  H(R,\Pi) \;=\; \tfrac12\,\Pi\cdot\Jm^{-1}\Pi ,
\end{equation}
which is left-invariant ($\mathbf{d}_RH = 0$; a spacecraft has no gravitational
potential in this model). Since
$\dpair{\ad^{*}_{\Omega}\Pi}{\zeta} = \Pi\cdot(\Omega\times\zeta) = (\Pi\times\Omega)\cdot\zeta$,
Lemma~\ref{lem:triv} gives $\ad^{*}_{\Omega}\Pi = \Pi\times\Omega$ and hence
\begin{equation}
  \boxed{\;\dot R = R\hat\Omega, \qquad \dot\Pi = \Pi\times\Omega + \tau
  \quad\Longleftrightarrow\quad
  \Jm\dot\Omega + \Omega\times\Jm\Omega = \tau . \;}
  \label{eq:euler}
\end{equation}
These are Euler's equations; the second is the Lie--Poisson equation on
$\so^{*}$ for the $(-)$ bracket, and equivalently the Euler--Poincar\'e equation
for the reduced Lagrangian $\ell(\Omega)=\tfrac12\Omega\cdot\Jm\Omega$. With
$\tau\equiv0$, Note~I, Prop.~3.1 applies verbatim: $\Lie_{X_H}\omega_0=0$ and the
drift is symplectic, hence (Note~I, Prop.~3.2) never contracting.

\subsection{Affine connection form}

\begin{lemma}
\label{lem:so3-conn}
For the left-invariant metric $\Jm$ on $\SO$,
\begin{equation}
  \Gamma(x,y) \;=\; \tfrac12\Big( x\times y + \Jm^{-1}\big(x\times\Jm y\big)
                    + \Jm^{-1}\big(y\times\Jm x\big)\Big),
  \qquad
  \Gamma(x,x) = \Jm^{-1}(x\times\Jm x),
  \label{eq:so3-Gamma}
\end{equation}
and consequently \eqref{eq:euler} is equivalent to
\begin{equation}
  \Dt\Omega \;=\; \dot\Omega + \Jm^{-1}\big(\Omega\times\Jm\Omega\big) \;=\; \Jm^{-1}\tau,
  \qquad\text{i.e.}\qquad
  \Gflat(\nabla_v v) = \tau .
  \label{eq:so3-accs}
\end{equation}
\end{lemma}

\begin{proof}
$\adt_x y = \Jm^{-1}(\Jm y\times x)$, since
$\Jm(\adt_xy)\cdot z = \Jm y\cdot(x\times z) = (\Jm y\times x)\cdot z$; substitute
into \eqref{eq:LC}. Setting $y=x$ gives $\Gamma(x,x) = -\adt_xx = \Jm^{-1}(x\times\Jm x)$.
Then \eqref{eq:euler} reads $\dot\Omega = \Jm^{-1}\tau - \Jm^{-1}(\Omega\times\Jm\Omega)$.
\end{proof}

Comparing with Note~II, Lemma~1.1: here $\mathcal{H}=\Jm$ is constant, so
$\dot{\mathcal{H}} - 2\mathcal{C} = -2\mathcal{C}$ and the skew-symmetry of the
gyroscopic matrix $\mathcal{C}(\Omega)W = \Jm\Gamma(\Omega,W)$ is the whole content
of $\nabla\Gm=0$ in this case.

\subsection{Tracking error system}

Let $R_d(t)$ satisfy $\dot R_d = R_d\hat\Omega_d$ with $\Omega_d,\dot\Omega_d$ known
and bounded. Define
\begin{equation}
  R_e \coloneqq R_d^{\top}R \in \SO,
  \qquad
  \Omr \coloneqq R_e^{\top}\Omega_d,
  \qquad
  \Ome \coloneqq \Omega - \Omr .
  \label{eq:so3-err}
\end{equation}

\begin{lemma}[The error system is again left-invariant]
\label{lem:err-kin}
$\dot R_e = R_e\hat\Ome$ and $\dot\Omr = R_e^{\top}\dot\Omega_d - \Ome\times\Omr$.
\end{lemma}

\begin{proof}
$\dot R_e = -\hat\Omega_dR_d^{\top}R + R_d^{\top}R\hat\Omega
= R_e\big(\hat\Omega - R_e^{\top}\hat\Omega_dR_e\big) = R_e(\hat\Omega - \widehat{R_e^{\top}\Omega_d})$,
using $R^{\top}\hat xR = \widehat{R^{\top}x}$. For the second,
$\dot\Omr = \dot R_e^{\top}\Omega_d + R_e^{\top}\dot\Omega_d$ and
$\dot R_e^{\top} = -\hat\Ome R_e^{\top}$.
\end{proof}

So $\Omr$ is the transport of the reference angular velocity to the current body
frame --- a transport map in the sense of Note~II, Def.~2.1 --- and $(R_e,\Ome)$ is
a curve on $\SO$ with body velocity $\Ome$. \emph{Tracking on $\SO$ is
stabilisation of a left-invariant system on $\SO$}, and all of Note~I applies on
the error copy.

%%%%%%%%%%%%%%%%%%%%%%%%%%%%%%%%%%%%%%%%%%%%%%%%%%%%%%%%%%%%%%%%%%%%%%%%%%%%%%%
\section{The shaping potential and its Hessian}
\label{sec:so3-pot}

\begin{definition}
For $\Km = \Km^{\top}\succ0$ define $\varphi\in\Cinf{\SO}$ and $\eR$ by
\begin{equation}
  \varphi(R_e) \;=\; \tfrac12\,\tr\!\big(\Km(I - R_e)\big),
  \qquad
  \eR \;=\; \tfrac12\big(\Km R_e - R_e^{\top}\Km\big)^{\vee}.
\end{equation}
\end{definition}

\begin{lemma}[Gradient and Hessian]
\label{lem:hess}
Let $\mathbf{d}\varphi\in\Rm^3$ be the left-trivialised differential, defined by
$\tfrac{d}{dt}\big|_0\varphi(R_e e^{t\hat\eta}) = \mathbf{d}\varphi\cdot\eta$.
Then $\mathbf{d}\varphi = \eR$, $\grad\varphi = \Jm^{-1}\eR$, and
$\Hess\varphi(\eta,\eta) = \eta\cdot\mathbf{H}(R_e)\,\eta$ with
\begin{equation}
  \boxed{\;
  \mathbf{H}(R_e) \;=\; \tfrac12\Big(\tr(\Km R_e)\,I - \sym\!\big(\Km R_e\big)\Big)
   \;-\; \sym\!\Big(\Jm\,\widehat{\Jm^{-1}\eR}\Big) .
  \;}
  \label{eq:Hmat}
\end{equation}
Hence $\Hess^{\sharp}\varphi = \Jm^{-1}\mathbf{H}$, whose $\Gm$-spectrum is the
generalised spectrum of the pencil $(\mathbf{H},\Jm)$. At $R_e = I$,
$\mathbf{H}(I) = \tfrac12\big(\tr(\Km)I - \Km\big)$.
\end{lemma}

\begin{proof}
Writing $A = \Km R_e$, $\tfrac{d}{dt}\big|_0\varphi(R_ee^{t\hat\eta}) = -\tfrac12\tr(A\hat\eta)$;
for any $A$, $\tr(A\hat\eta) = \tr(\hat a\hat\eta) = -2a\cdot\eta$ with
$a=\tfrac12(A-A^{\top})^{\vee}$, giving $\mathbf{d}\varphi=\eR$. Differentiating
$\eR$ again along $R_ee^{t\hat\eta}$ uses the identity
\begin{equation}
  \big(A\hat\eta + \hat\eta A^{\top}\big)^{\vee} \;=\; \big(\tr(A)I - A^{\top}\big)\eta
  \qquad\text{for every }A\in\Rm^{3\times3},
  \label{eq:vee-id}
\end{equation}
so $\tfrac{d^2}{dt^2}\big|_0\varphi(R_ee^{t\hat\eta}) = \eta\cdot E\eta$ with
$E = \tfrac12(\tr(\Km R_e)I - R_e^{\top}\Km)$ and
$\sym(E) = \tfrac12(\tr(\Km R_e)I - \sym(\Km R_e))$. The curve
$t\mapsto R_ee^{t\hat\eta}$ is \emph{not} a geodesic (one-parameter subgroups are
geodesics only for bi-invariant metrics), and has covariant acceleration
$\Gamma(\eta,\eta)$; therefore
$\Hess\varphi(\eta,\eta) = \tfrac{d^2}{dt^2}\big|_0\varphi(R_ee^{t\hat\eta}) - \mathbf{d}\varphi\cdot\Gamma(\eta,\eta)$.
Finally, with $w=\Jm^{-1}\eR$,
$\eR\cdot\Gamma(\eta,\eta) = w\cdot(\eta\times\Jm\eta) = \eta^{\top}\Jm\hat w\,\eta
= \eta^{\top}\sym(\Jm\hat w)\eta$.
\end{proof}

\begin{remark}
The second term of \eqref{eq:Hmat} vanishes identically when $\Jm = jI$, since
$\sym(\hat w)=0$. It is the correction that distinguishes a left-invariant from a
bi-invariant metric, and it is precisely what a Euclidean derivation of the
Hessian would miss.
\end{remark}

\begin{proposition}[Isotropising gain, and the inertia triangle inequality]
\label{prop:isotropise}
$\Hess^{\sharp}\varphi(I) = k\,\id$ if and only if
\begin{equation}
  \Km \;=\; k\big(\tr(\Jm)\,I - 2\Jm\big),
  \qquad k>0 .
  \label{eq:iso-gain}
\end{equation}
This $\Km$ is positive-definite if and only if
$\lambda_{\max}(\Jm) < \tr(\Jm) - \lambda_{\max}(\Jm)$, i.e.\ iff the largest
principal moment of inertia is strictly less than the sum of the other two ---
which holds for every rigid body that is not a planar lamina.
\end{proposition}

\begin{proof}
$\Jm^{-1}\mathbf{H}(I) = k\id \iff \tfrac12(\tr(\Km)I-\Km) = k\Jm$. Taking traces,
$\tr\Km = k\tr\Jm$; substituting gives \eqref{eq:iso-gain}, and conversely
$\tr\big(k(\tr(\Jm)I-2\Jm)\big) = k\tr\Jm$ returns $\mathbf{H}(I) = k\Jm$. Positivity
of $\Km$ is $\tr(\Jm) - 2\lambda_{\max}(\Jm)>0$. Writing $\Jm$'s eigenvalues as the
principal moments $J_1\le J_2\le J_3$, this is $J_1+J_2>J_3$, the standard
triangle inequality, with equality exactly for a lamina.
\end{proof}

Under \eqref{eq:iso-gain} the shaped stiffness is isotropic \emph{at} the target,
so by Note~I, Remark~6.5 arbitrarily small damping is admissible in the
infinitesimal limit; the finite-region result below quantifies the price of a
finite contraction region.

The critical points of $\varphi$ are $R_e=I$ and the three rotations by $\pi$ about
the eigenaxes of $\Km$, with values $\tr(\Km)-\kappa_i$ where $\kappa_i$ are the
eigenvalues of $\Km$. Hence
\begin{equation}
  \varphi(R_e) < \tr(\Km) - \lambda_{\max}(\Km)
  \quad\Longrightarrow\quad
  R_e = I \text{ is the only critical point in that sublevel set}.
  \label{eq:crit}
\end{equation}

%%%%%%%%%%%%%%%%%%%%%%%%%%%%%%%%%%%%%%%%%%%%%%%%%%%%%%%%%%%%%%%%%%%%%%%%%%%%%%%
\section{Curvature of $(\SO,\Gm_{\Jm})$}
\label{sec:so3-curv}

Note~I, Thm.~6.4 needs the Jacobi operator $\Jac_v(u)=R(u,v)v$. From
Lemma~\ref{lem:LC}, for left-invariant fields
\begin{equation}
  R(x,y)z \;=\; \Gamma\big(x,\Gamma(y,z)\big) - \Gamma\big(y,\Gamma(x,z)\big)
                - \Gamma\big(x\times y,\;z\big).
  \label{eq:curv}
\end{equation}

\begin{lemma}
\label{lem:curv}
\begin{enumerate}[label=(\roman*),leftmargin=2.1em,itemsep=2pt,topsep=2pt]
\item If $\Jm = jI$ then $\Gamma(x,y)=\tfrac12 x\times y$ and $(\SO,\Gm_{\Jm})$ has
      constant sectional curvature $\Sect \equiv 1/(4j) > 0$.
\item For general $\Jm$, define
      \begin{equation}
        C_R \;\coloneqq\; \sup_{\normJ{v}=1}\; \big\|\,\Jac_v\,\big\|_{\Jm}
        \qquad\text{so that}\qquad
        \big\|\Jac_v\big\|_{\Jm} \le C_R\,\normJ{v}^2 .
        \label{eq:CR}
      \end{equation}
      $\Jac_v$ is $\Gm$-symmetric, and $C_R$ is a two-parameter supremum over the
      unit $\Jm$-sphere, computable to machine precision. A crude closed-form
      bound is $C_R \le 2c_\Gamma^2 + c_\Gamma c_{\ad}$ with
      $c_{\ad} = \sqrt{\bar\jmath}/\underline\jmath$ and
      $c_\Gamma = \tfrac12\big(\sqrt{\bar\jmath}/\underline\jmath + 2\bar\jmath/\underline\jmath^{3/2}\big)$,
      $\underline\jmath = \lambda_{\min}(\Jm)$, $\bar\jmath = \lambda_{\max}(\Jm)$.
\item If $\Jm \ne jI$, some sectional curvatures are negative.
\end{enumerate}
\end{lemma}

Item (iii) is Milnor's observation for left-invariant metrics; for
$\Jm=\mathrm{diag}(1,1,8)$ the sectional curvatures range over
$[-5.00,\,2.00]$ numerically. This matters: by Note~I, \S5, negative curvature
subtracts from the effective stiffness at rate $\normJ{\Ome}^2$ and \emph{cannot be
repaired by damping}, because $\Damp$ acts only on the fibre block. An inertially
asymmetric spacecraft therefore has an intrinsically velocity-bounded contraction
region. The closed-form bound in (ii) is very conservative --- for
$\Jm = \begin{bsmallmatrix}2&.3&.1\\.3&3.5&.2\\.1&.2&5\end{bsmallmatrix}$ it gives
$14.7$ against a true $C_R = 0.176$ --- so \eqref{eq:CR} should be evaluated
directly.

%%%%%%%%%%%%%%%%%%%%%%%%%%%%%%%%%%%%%%%%%%%%%%%%%%%%%%%%%%%%%%%%%%%%%%%%%%%%%%%
\section{Controller, metric, and contraction theorem}
\label{sec:so3-main}

\subsection{Controller}

\begin{definition}[Attitude tracking law]
\label{def:so3-ctrl}
With $\Km$ as in \eqref{eq:iso-gain}, damping gain $d>0$, and $R_e,\Omr,\Ome$ from
\eqref{eq:so3-err}, set
\begin{equation}
  \boxed{\;
  \tau \;=\; \Jm\big(R_e^{\top}\dot\Omega_d - \Ome\times\Omr\big)
   \;+\; \Ome\times\Jm\Omr + \Omr\times\Jm\Ome + \Omr\times\Jm\Omr
   \;-\; \eR \;-\; d\,\Jm\Ome .
  \;}
  \label{eq:so3-tau}
\end{equation}
\end{definition}

The first line is feedforward ($\Jm\dot\Omr$ plus the gyroscopic terms generated by
the reference); the second is potential shaping plus damping. For a setpoint
($\Omega_d\equiv0$) it collapses to $\tau = -\eR - d\,\Jm\Omega$, the classical PD
law of \cite[\S11.1]{BL04}.

\begin{proposition}[Closed loop]
\label{prop:so3-cl}
Under \eqref{eq:so3-tau}, the error system is exactly Note~I, eq.~(19) on the error
copy of $\SO$ with $V\equiv0$, $\Damp = d\,\id$:
\begin{equation}
  \dot R_e = R_e\hat\Ome,
  \qquad
  \Dt\Ome \;=\; \dot\Ome + \Jm^{-1}\big(\Ome\times\Jm\Ome\big)
             \;=\; -\grad\varphi(R_e) - d\,\Ome .
  \label{eq:so3-cl}
\end{equation}
\end{proposition}

\begin{proof}
\eqref{eq:so3-cl} requires
$\Jm\dot\Ome = -\Ome\times\Jm\Ome - \eR - d\Jm\Ome$. Since
$\dot\Ome = \dot\Omega - \dot\Omr$ and $\Jm\dot\Omega = \tau - \Omega\times\Jm\Omega$,
this holds iff
$\tau = \Jm\dot\Omr + \big(\Omega\times\Jm\Omega - \Ome\times\Jm\Ome\big) - \eR - d\Jm\Ome$.
Expanding $\Omega = \Ome+\Omr$ gives
$\Omega\times\Jm\Omega - \Ome\times\Jm\Ome = \Ome\times\Jm\Omr + \Omr\times\Jm\Ome + \Omr\times\Jm\Omr$,
and $\Jm\dot\Omr = \Jm(R_e^{\top}\dot\Omega_d - \Ome\times\Omr)$ by
Lemma~\ref{lem:err-kin}.
\end{proof}

\subsection{Contraction region and the constants $\mu,\sigma$}

Let $\mathcal{E}(R_e,\Ome) = \tfrac12\normJ{\Ome}^2 + \varphi(R_e)$. Along
\eqref{eq:so3-cl}, $\dot{\mathcal{E}} = -d\normJ{\Ome}^2 \le 0$, so sublevel sets
are forward invariant. Fix $c_0>0$ and put
\begin{equation}
  \mathcal{W} \;\coloneqq\; \big\{(R_e,\Ome) \;:\; \mathcal{E}(R_e,\Ome) < c_0 \big\}
  \;\subset\; T^{*}\SO ,
  \label{eq:W}
\end{equation}
identified with a subset of $T^{*}\SO$ by $\Pi_e = \Jm\Ome$. Define, with
$\Jm = L L^{\top}$ a Cholesky factor,
\begin{equation}
  \mu \;\coloneqq\; \inf_{\mathcal{W}} \lambda_{\min}\!\big(\sym\Psi\big),
  \qquad
  \sigma \;\coloneqq\; \sup_{\mathcal{W}} \big\|\Psi\big\|_2,
  \qquad
  \Psi \;\coloneqq\; L^{-1}\mathbf{H}(R_e)L^{-\top} + L^{\top}\Jac_{\Ome}L^{-\top},
  \label{eq:musigma}
\end{equation}
i.e.\ $\Psi$ is the $\Jm$-orthonormal representation of
$\Hess^{\sharp}\varphi + \Jac_{\Ome}$. Both are explicit given $\Jm,\Km,c_0$ and
are bounded by
\begin{equation}
  \mu \;\ge\; \inf_{\varphi(R_e)<c_0}\lambda_{\min}\big(L^{-1}\mathbf{H}L^{-\top}\big) - 2c_0C_R,
  \qquad
  \sigma \;\le\; \sup_{\varphi(R_e)<c_0}\big\|L^{-1}\mathbf{H}L^{-\top}\big\|_2 + 2c_0C_R ,
  \label{eq:musigma-bnd}
\end{equation}
using $\normJ{\Ome}^2 < 2c_0$ on $\mathcal{W}$ and \eqref{eq:CR}.

\subsection{Contraction metric}

\begin{definition}[Contraction metric on $T^{*}\SO$]
\label{def:so3-metric}
Let $w\in T_\alpha(T^{*}\SO)$ split as $(u,\xi)$ in the $\nabla$-induced splitting
of Note~I, \S4.1, with $u,\xi\in\Rm^3$ the left-trivialised (body) representations
of the horizontal and vertical components. With $\mu,\sigma$ from
\eqref{eq:musigma} and $d>0$, set
\begin{equation}
  \boxed{\;
  G_\alpha(w,w) \;=\; a\;u\cdot\Jm u \;+\; 2b\;u\cdot\Jm\xi \;+\; \xi\cdot\Jm\xi,
  \qquad
  b = \tfrac{d}{2},
  \quad
  a = \tfrac{d^2}{2} + \tfrac{\mu+\sigma}{2}. \;}
  \label{eq:so3-G}
\end{equation}
\end{definition}

This is Note~I, Def.~6.1 and eq.~(23) with $c=1$: horizontal block $a\,\pi^{*}\Gm$,
vertical block $\Gm^{-1}$, and the base--fibre cross term $b\,\Pi$ that Note~I,
Prop.~5.1 shows is indispensable. Positive-definiteness is
$a > b^2$, i.e.\ $d^2/2 + (\mu+\sigma)/2 > d^2/4$, automatic.

\subsection{Main theorem}

\begin{theorem}[Contraction of the attitude tracking loop]
\label{thm:so3}
Let $\Jm\succ0$, let $\Km$ satisfy \eqref{eq:iso-gain} (or be any $\Km\succ0$), let
$c_0>0$ satisfy \eqref{eq:crit} and be small enough that $\mu>0$ in
\eqref{eq:musigma}, and let $\mathcal{W}$ be $K$-reachable in the metric
\eqref{eq:so3-G}. If the damping gain obeys
\begin{equation}
  \boxed{\;\; d \;>\; \frac{\sigma-\mu}{2\sqrt{\mu}} \;\;}
  \label{eq:so3-cond}
\end{equation}
then the closed loop \eqref{eq:so3-cl} is contracting on $\mathcal{W}$ with respect
to $G$ of \eqref{eq:so3-G}, with rate
\begin{equation}
  \lambda = \frac{\lambda_{\min}(\mathcal{Q})}{\lambda_{\max}(\mathcal{P})},
  \quad
  \mathcal{Q} = \begin{bmatrix}\tfrac{\mu d}{2} & -\tfrac{\kappa}{2}\\[2pt]
                               -\tfrac{\kappa}{2} & \tfrac{d}{2}\end{bmatrix},
  \quad
  \mathcal{P} = \begin{bmatrix}a & b\\ b& 1\end{bmatrix},
  \quad
  \kappa \le \tfrac{\sigma-\mu}{2}.
\end{equation}
Consequently any two error trajectories in $\mathcal{W}$ satisfy
$d_G(\Phi_t(\alpha_0),\Phi_t(\alpha_1)) \le Ke^{-\lambda t}d_G(\alpha_0,\alpha_1)$,
and since $(R_e,\Ome)=(I,0)$ is an equilibrium of \eqref{eq:so3-cl} lying in
$\mathcal{W}$, every trajectory in $\mathcal{W}$ converges to it exponentially:
$R\to R_d$ and $\Omega\to\Omr$.
\end{theorem}

\begin{proof}
By Prop.~\ref{prop:so3-cl} the closed loop is Note~I, eq.~(19) on $Q=\SO$ with
$V\equiv0$, shaping potential $\varphi$, and $\Damp = d\,\id$, which is
$\nabla$-parallel (Note~I, Assumption~6.3). Full actuation
(Note~I, Assumption~6.2) holds because $\tau$ ranges over all of $\so^{*}\cong\Rm^3$.
Note~I, eq.~(20) then gives the covariant variational system
\begin{equation}
  \Dt u = \xi,
  \qquad
  \Dt\xi = -\Psi_\alpha(u) - d\,\xi,
  \qquad
  \Psi_\alpha = \Hess^{\sharp}\varphi(R_e) + \Jac_{\Ome},
  \label{eq:so3-var}
\end{equation}
with $\Hess^{\sharp}\varphi = \Jm^{-1}\mathbf{H}$ from Lemma~\ref{lem:hess} and
$\Jac_{\Ome}$ from \eqref{eq:curv}. Both are $\Gm$-symmetric, and by
\eqref{eq:musigma} $\spec_{\Gm}(\Psi_\alpha)\subset[\mu,\sigma]$ uniformly on
$\mathcal{W}$. $\mathcal{W}$ is forward invariant since
$\dot{\mathcal{E}} = -d\normJ{\Ome}^2\le0$, and the closed loop is forward complete
on it because $\overline{\mathcal{W}}$ is compact ($\SO$ compact, $\Ome$ bounded).
All hypotheses of Note~I, Thm.~6.4 are met, and \eqref{eq:so3-cond} is its
condition (22) with $c=1$. The metric \eqref{eq:so3-G} is its eq.~(23).

For completeness, the one-line verification: writing
$G(w,w) = a\normJ{u}^2 + 2b\ipJ{u}{\xi} + \normJ{\xi}^2$ and using
$\nabla\Gm=0$ together with \eqref{eq:so3-var},
\[
  \tfrac12\tfrac{d}{dt}G(w,w)
  = -b\ipJ{\Psi u}{u} - (d-b)\normJ{\xi}^2
    + \ipJ{\big(a\,\id - bd\,\id - \Psi\big)\xi}{u} .
\]
The choice $a = bd + \tfrac{\mu+\sigma}{2}$ centres the residual operator so that
$\kappa \coloneqq \sup_{\mathcal{W}}\norm{a\id-bd\id-\Psi} \le \tfrac{\sigma-\mu}{2}$,
and $b=d/2$ maximises $b\mu(d-b)$, reducing $\mathcal{Q}\succ0$ to
\eqref{eq:so3-cond}. Then
$\tfrac{d}{dt}G(w,w) \le -2\lambda G(w,w)$, and \cite[Thm.~2.3]{SB13} gives the
distance bound.
\end{proof}

\begin{remark}[Numerical verification]
\label{rem:numeric}
For $\Jm = \begin{bsmallmatrix}2&0.3&0.1\\0.3&3.5&0.2\\0.1&0.2&5\end{bsmallmatrix}$,
$\Km = 0.7(\tr(\Jm)I-2\Jm)$ (eigenvalues $0.304,2.413,4.633 \succ 0$, consistent with
Prop.~\ref{prop:isotropise}), $d=1.2$ and $c_0 = 0.3984$, direct evaluation of
\eqref{eq:musigma} gives $\mu = 0.5485$, $\sigma = 0.8272$, so
\eqref{eq:so3-cond} requires $d > 0.1882$ --- satisfied. The metric is
$(a,b,c) = (1.4079,\,0.6,\,1)$ with $ac-b^2 = 1.048>0$ and $\lambda = 0.1699$.
Integrating \eqref{eq:so3-var} along a closed-loop trajectory from
$\mathcal{E}=c_0$, the observed $\sup_t \dot G/G$ is $-1.011$, comfortably below
the certified $-2\lambda = -0.340$: the certificate is valid and, as expected from
the Young-inequality step, conservative.
\end{remark}

\begin{remark}[Partial-contraction alternative]
Note~II, Thm.~3.4 gives a cheaper claim: with
$\tau = \Jm\dot\Omr + \mathcal{C}(\Omega)\Omr - \Jm\Kgain\Ome$,
$\mathcal{C}(\Omega)W = \Jm\Gamma(\Omega,W)$, the virtual system on
$\gamma^{*}T\SO$ is contracting in $\Gm=\Jm$ alone, with rate
$\lambda_{\min}(\sym_{\Jm}\Kgain)$, giving $\Omega\to\Omr$ but not $R\to R_d$.
That argument needs no Hessian and no curvature --- precisely because, as Note~II,
\S5 explains, the virtual system suppresses the horizontal channel where
$\Hess^{\sharp}\varphi$ and $\Jac_v$ live. Theorem~\ref{thm:so3} is the stronger,
more expensive claim.
\end{remark}

\begin{remark}[Global obstruction]
$T^{*}\SO$ deformation-retracts onto $\SO\simeq\Rm P^3$, which is not
contractible, so by Note~I, \S5 (following \cite[\S2]{SB13} and Bhat--Bernstein)
$\mathcal{W}$ can never be all of $T^{*}\SO$. Condition \eqref{eq:crit} is the
concrete form this takes: the three $\pi$-rotations about the eigenaxes of $\Km$
are the unavoidable non-target critical points.
\end{remark}

%%%%%%%%%%%%%%%%%%%%%%%%%%%%%%%%%%%%%%%%%%%%%%%%%%%%%%%%%%%%%%%%%%%%%%%%%%%%%%%
\part*{Part B. Quadrotor on $\SE$}
\addcontentsline{toc}{section}{Part B. Quadrotor on $\SE$}
\setcounter{section}{0}
\renewcommand{\thesection}{B.\arabic{section}}
\setcounter{theorem}{0}
\renewcommand{\thetheorem}{B.\arabic{section}.\arabic{theorem}}

\section{Dynamics}
\label{sec:se3-dyn}

\subsection{Configuration manifold, metric, potential}

$Q = \SE$, $q = (x,R)$ with $x\in\Rm^3$ the centre of mass in the inertial frame.
$\se \cong \Rm^3\times\Rm^3 \ni \xi = (V,\Omega)$ with
\begin{equation}
  \big[(V_1,\Omega_1),(V_2,\Omega_2)\big]
  \;=\; \big(\Omega_1\times V_2 - \Omega_2\times V_1,\; \Omega_1\times\Omega_2\big).
\end{equation}
The left-invariant frame gives $\dot x = RV$, $\dot R = R\hat\Omega$. The kinetic
energy metric is left-invariant with generalised inertia
\begin{equation}
  \Ig = \begin{bmatrix} mI_3 & 0\\ 0 & \Jm\end{bmatrix},
  \qquad
  \tfrac12\ip{\xi}{\xi}_{\Ig} = \tfrac12 m\norm{V}^2 + \tfrac12\Omega\cdot\Jm\Omega .
\end{equation}
Gravity is $V_g(x,R) = mg\,e_3\cdot x$; its left-trivialised differential is
$\mathbf{d}V_g = (mg\,\Gamma_{\!3},\,0)$ with the \emph{advected direction}
$\Gamma_{\!3}\coloneqq R^{\top}e_3$, so $\grad V_g = (g\Gamma_{\!3},0)$. The potential is
\emph{not} left-invariant, which is why only a semidirect-product reduction is
available.

\subsection{Hamiltonian, Lie--Poisson, Euler--Poincar\'e}

With $\mu = \Ig\xi = (mV,\Jm\Omega)$ and
$H = \tfrac{1}{2m}\norm{mV}^2 + \tfrac12\Omega\cdot\Jm\Omega + mge_3\cdot x$,
one computes
$\ad^{*}_{(V,\Omega)}(p,\Pi) = (p\times\Omega,\; p\times V + \Pi\times\Omega)$,
which at $p=mV$ reduces to $(mV\times\Omega,\;\Jm\Omega\times\Omega)$.
Lemma~\ref{lem:triv} with body wrench $F = (fe_3,\,M)$ then gives
\begin{equation}
  \boxed{\;
  \begin{aligned}
    m\dot V &= mV\times\Omega + f e_3 - mg\,\Gamma_{\!3}, \\
    \Jm\dot\Omega &= \Jm\Omega\times\Omega + M, \\
    \dot\Gamma_{\!3} &= \Gamma_{\!3}\times\Omega,
    \qquad \dot x = RV, \qquad \dot R = R\hat\Omega .
  \end{aligned}\;}
  \label{eq:se3-EP}
\end{equation}
The first three lines are the Euler--Poincar\'e equations on
$\se^{*}\ltimes\Rm^3$; eliminating $V$ in favour of $\dot x$ returns the familiar
inertial-frame model
\begin{equation}
  m\ddot x = f R e_3 - mg\,e_3,
  \qquad
  \Jm\dot\Omega + \Omega\times\Jm\Omega = M .
  \label{eq:se3-NE}
\end{equation}
By \eqref{eq:covaccel} this is $\Gflat(\nabla_vv) = -dV_g + F$, i.e.\ Note~I,
eq.~(1); with $F\equiv0$ the drift is symplectic (Note~I, Prop.~3.1).

\subsection{Inputs and the under-actuation obstruction}

Four rotor thrusts $f_i\ge0$ map to $(f,M)$ through the constant, invertible mixer
determined by arm length and torque coefficient, so $(f,M)\in\Rm\times\Rm^3$ is an
equivalent input --- subject to $f_i\in[0,f_{\max}]$, which we return to in
\S\ref{sec:se3-caveats}.

\begin{proposition}[Note~I, Assumption~6.2 fails]
\label{prop:underact}
The input covector fields span
$\mathcal{F}_q = \{(ce_3,\,M) : c\in\Rm,\ M\in\Rm^3\} \subset \se^{*}$, of dimension
$4 < 6 = \dim\se^{*}$. The unactuated directions are the body $e_1,e_2$
translations. Hence no choice of $(f,M)$ can realise the shaping-plus-damping law
of Note~I, eq.~(18), and Note~I, Thm.~6.4 does not apply.
\end{proposition}

This is not a technicality that can be patched by a different metric: Note~I,
\S8 flags full actuation as used essentially, and the missing directions are
exactly the ones a quadrotor must generate by tilting. What follows is the
standard resolution --- generate the missing translational force by attitude ---
recast so that the coupling is handled by \cite[Lem.~3.2]{SB13}.

%%%%%%%%%%%%%%%%%%%%%%%%%%%%%%%%%%%%%%%%%%%%%%%%%%%%%%%%%%%%%%%%%%%%%%%%%%%%%%%
\section{Decomposition into two contracting subsystems}
\label{sec:se3-decomp}

\subsection{Virtual translational input}

Write $b_3 \coloneqq Re_3$ and define the \emph{virtual translational input}
$\upsilon \coloneqq f b_3 \in \Rm^3$. If $\upsilon$ were free, the translational
subsystem $m\ddot x = \upsilon - mge_3$ on $Q_{\mathrm{tr}} = \Rm^3$ would be fully
actuated. It is not: $\upsilon$ is constrained to the ray through $b_3$.

Given a reference $x_d(t)$, set $e_x = x-x_d$, $e_v = \dot x - \dot x_d$, gains
$k_x,k_v>0$, and
\begin{equation}
  F_{\mathrm{des}} \;\coloneqq\; -k_xe_x - k_ve_v + mg\,e_3 + m\ddot x_d,
  \qquad
  f \;\coloneqq\; F_{\mathrm{des}}\cdot b_3,
  \qquad
  b_{3c} \;\coloneqq\; \frac{F_{\mathrm{des}}}{\norm{F_{\mathrm{des}}}} .
  \label{eq:Fdes}
\end{equation}
The commanded attitude $R_c$ has third column $b_{3c}$, first column
$b_{1c} = \Pi_{b_{3c}^{\perp}}(b_{1d})/\norm{\cdot}$ for a desired heading $b_{1d}$
with $b_{1d}\not\parallel b_{3c}$, and $b_{2c}=b_{3c}\times b_{1c}$.

\begin{lemma}[Translational error dynamics]
\label{lem:tr-err}
With \eqref{eq:Fdes},
\begin{equation}
  m\,\ddot e_x + k_v e_v + k_x e_x \;=\; \Delta,
  \qquad
  \Delta \coloneqq -\big(I - b_3b_3^{\top}\big)F_{\mathrm{des}},
  \qquad
  \norm{\Delta} = \norm{F_{\mathrm{des}}}\,\sin\theta ,
  \label{eq:tr-err}
\end{equation}
where $\theta = \angle(b_3,b_{3c})$. In particular $\Delta = 0$ iff $b_3 = \pm b_{3c}$.
\end{lemma}

\begin{proof}
$m\ddot e_x = fb_3 - mge_3 - m\ddot x_d = (F_{\mathrm{des}}\cdot b_3)b_3 - F_{\mathrm{des}} - k_xe_x - k_ve_v$
after substituting $mge_3+m\ddot x_d = F_{\mathrm{des}}+k_xe_x+k_ve_v$; and
$(F_{\mathrm{des}}\cdot b_3)b_3 - F_{\mathrm{des}} = -(I-b_3b_3^{\top})F_{\mathrm{des}}$.
Finally $(I-b_3b_3^{\top})F_{\mathrm{des}} = \norm{F_{\mathrm{des}}}(I-b_3b_3^{\top})b_{3c}$
and $\norm{(I-b_3b_3^{\top})b_{3c}} = \sin\theta$.
\end{proof}

So $\Delta$ is exactly the component of the desired force that the current attitude
cannot deliver, and it is the interconnection signal.

\subsection{The translational subsystem is contracting for every $k_v>0$}

\begin{proposition}[Translational certificate]
\label{prop:tr}
Consider \eqref{eq:tr-err} with $\Delta\equiv0$ as a mechanical system on
$Q_{\mathrm{tr}}=\Rm^3$ with $\Gm_{\mathrm{tr}} = mI$, shaping potential
$\varphi_{\mathrm{tr}} = \tfrac{k_x}{2}\norm{e_x}^2$ and damping
$\Damp_{\mathrm{tr}} = \tfrac{k_v}{m}\id$. Then $Q_{\mathrm{tr}}$ is flat, so
$\Jac_v \equiv 0$ and
\begin{equation}
  \Psi_{\mathrm{tr}} = \Hess^{\sharp}\varphi_{\mathrm{tr}} = \tfrac{k_x}{m}\,\id,
  \qquad\text{hence}\qquad
  \mu_{\mathrm{tr}} = \sigma_{\mathrm{tr}} = \tfrac{k_x}{m}.
\end{equation}
Condition \eqref{eq:so3-cond} reads $d_{\mathrm{tr}} = k_v/m > 0$: \emph{every}
positive damping gain works. The certificate is
\begin{equation}
  G_{\mathrm{tr}}(w,w) \;=\; m\Big( a_{\mathrm{tr}}\norm{u}^2 + 2 b_{\mathrm{tr}}\,u\cdot\xi + \norm{\xi}^2 \Big),
  \qquad
  b_{\mathrm{tr}} = \frac{k_v}{2m},
  \quad
  a_{\mathrm{tr}} = \frac{k_v^2}{2m^2} + \frac{k_x}{m},
  \label{eq:Gtr}
\end{equation}
with rate $\lambda_{\mathrm{tr}} = \lambda_{\min}(\mathcal{Q}_{\mathrm{tr}})/\lambda_{\max}(\mathcal{P}_{\mathrm{tr}})$,
$\mathcal{Q}_{\mathrm{tr}} = \mathrm{diag}\big(\tfrac{k_xk_v}{2m^2},\tfrac{k_v}{2m}\big)$,
$\mathcal{P}_{\mathrm{tr}} = \begin{bsmallmatrix}a_{\mathrm{tr}}&b_{\mathrm{tr}}\\ b_{\mathrm{tr}}&1\end{bsmallmatrix}$,
on any $\mathcal{W}_{\mathrm{tr}}$ that is a sublevel set of
$\tfrac{m}{2}\norm{e_v}^2 + \varphi_{\mathrm{tr}}$.
\end{proposition}

\begin{proof}
Note~I, Thm.~6.4 with $Q=\Rm^3$, whose curvature vanishes, so
$\Psi = \Hess^{\sharp}\varphi_{\mathrm{tr}}$ is a constant multiple of the identity
and $\kappa = 0$; this is Note~I, Remark~6.5.
\end{proof}

Note that the \emph{cross term} $2b_{\mathrm{tr}}u\cdot\xi$ is still essential:
Note~I, Prop.~5.1 rules out the block-diagonal energy metric here as everywhere.

\subsection{The attitude subsystem}

Apply Part~A with $R_d$ replaced by the commanded $R_c(t)$ of \eqref{eq:Fdes}, the
control being $M$ in place of $\tau$. This requires $R_c$ to be twice
differentiable with bounded $\Omega_c,\dot\Omega_c$, which holds on any region
where $F_{\mathrm{des}}$ is bounded away from zero and $x_d\in C^4$. Theorem~\ref{thm:so3}
then gives a contracting attitude subsystem on $\mathcal{W}_{\mathrm{att}}$ with
metric $G_{\mathrm{att}}$ of \eqref{eq:so3-G} and rate $\lambda_{\mathrm{att}}$.

\subsection{This is a feedback interconnection, not a cascade}

Two couplings are present:
\begin{enumerate}[label=(\roman*),leftmargin=2.1em,itemsep=2pt,topsep=2pt]
\item \textbf{attitude $\to$ translation}: through $\Delta$ in \eqref{eq:tr-err},
      driven by the attitude error;
\item \textbf{translation $\to$ attitude}: through $R_c$, which depends on
      $F_{\mathrm{des}}$ and hence on $(e_x,e_v)$.
\end{enumerate}
Because of (ii) the composite is \emph{not} a cascade, so \cite[Lem.~3.1]{SB13}
does not apply and the small-gain result \cite[Lem.~3.2]{SB13} is required.

%%%%%%%%%%%%%%%%%%%%%%%%%%%%%%%%%%%%%%%%%%%%%%%%%%%%%%%%%%%%%%%%%%%%%%%%%%%%%%%
\section{Interconnection theorem}
\label{sec:se3-main}

Following \cite[eq.~(11)]{SB13}, define the induced input--output gains of the two
contracting-with-inputs subsystems, with $h_1$ the output map of the translational
subsystem ($h_1(e_x,e_v) = F_{\mathrm{des}}$, which determines $R_c$) and $h_2$ that
of the attitude subsystem ($h_2(R_e,\Ome) = b_3$, which determines $\Delta$):
\begin{equation}
  \gamma_{\mathrm{tr}} \coloneqq
  \sup_{\mathcal{W}_{\mathrm{tr}}}\;\sup_{w\in T\mathcal{W}_{\mathrm{att}}}
  \frac{\tfrac12\big\|\partial_{\upsilon}X_{\mathrm{tr}}\circ (Th_2)\,w\big\|_{G_{\mathrm{tr}}}}
       {\norm{w}_{G_{\mathrm{att}}}},
  \qquad
  \gamma_{\mathrm{att}} \coloneqq
  \sup_{\mathcal{W}_{\mathrm{att}}}\;\sup_{v\in T\mathcal{W}_{\mathrm{tr}}}
  \frac{\tfrac12\big\|\partial_{R_c}X_{\mathrm{att}}\circ (Th_1)\,v\big\|_{G_{\mathrm{att}}}}
       {\norm{v}_{G_{\mathrm{tr}}}} .
  \label{eq:gains}
\end{equation}

\begin{theorem}[Contraction of the quadrotor tracking loop]
\label{thm:se3}
Let $\mathcal{W}_{\mathrm{tr}}\subset T^{*}\Rm^3$ and $\mathcal{W}_{\mathrm{att}}\subset T^{*}\SO$
be compact, forward-invariant, $K$-reachable regions on which
\begin{enumerate}[label=(H\arabic*),leftmargin=2.8em,itemsep=2pt,topsep=2pt]
\item \label{H1} the attitude loop satisfies Theorem~\ref{thm:so3} with rate
      $\lambda_{\mathrm{att}}$ and the translational loop satisfies
      Prop.~\ref{prop:tr} with rate $\lambda_{\mathrm{tr}}$;
\item \label{H2} $\norm{F_{\mathrm{des}}} \ge F_{\min} > 0$ and
      $\theta = \angle(b_3,b_{3c}) \le \theta_{\max} < \pi/2$, so the commanded
      thrust $f = \norm{F_{\mathrm{des}}}\cos\theta \ge F_{\min}\cos\theta_{\max} > 0$
      and $R_c$ is well defined and $C^2$;
\item \label{H3} the gains \eqref{eq:gains} are finite and satisfy the small-gain
      condition
      \begin{equation}
        \boxed{\;\;\gamma_{\mathrm{tr}}\,\gamma_{\mathrm{att}}
        \;<\; \lambda_{\mathrm{tr}}\,\lambda_{\mathrm{att}} .\;\;}
        \label{eq:smallgain}
      \end{equation}
\end{enumerate}
Then there exist $\alpha_1,\alpha_2>0$ such that the block-diagonal metric
\begin{equation}
  G \;=\; \alpha_1\,G_{\mathrm{tr}} \;\oplus\; \alpha_2\,G_{\mathrm{att}},
  \qquad
  \frac{\alpha_1}{\alpha_2} = \frac{\gamma_{\mathrm{att}}}{\gamma_{\mathrm{tr}}},
  \label{eq:se3-G}
\end{equation}
makes the closed loop contracting on
$\mathcal{W}_{\mathrm{tr}}\times\mathcal{W}_{\mathrm{att}}$, and $x\to x_d$,
$R\to R_c$ exponentially in $d_G$.
\end{theorem}

\begin{proof}
\ref{H1} supplies two contracting systems with inputs in the sense of
\cite[\S3]{SB13}: the translational one with input $\upsilon$ (equivalently
$\Delta$, by Lemma~\ref{lem:tr-err}) and metric \eqref{eq:Gtr}, the attitude one
with input $R_c$ and metric \eqref{eq:so3-G}. \ref{H2} guarantees both output maps
are smooth on the regions and that $f>0$, so the mixer inversion is well posed.
\ref{H3} is exactly hypothesis (12) of \cite[Lem.~3.2]{SB13}. That lemma's proof
computes, for $w = (v_1,w_2)$ in the product tangent space,
\[
  \ip{\overset{G}{\nabla}_{w}X}{w}_{G}
  \le -\alpha_1\lambda_{\mathrm{tr}}\norm{v_1}^2_{G_{\mathrm{tr}}}
      -\alpha_2\lambda_{\mathrm{att}}\norm{w_2}^2_{G_{\mathrm{att}}}
      +\big(\alpha_1\gamma_{\mathrm{tr}}+\alpha_2\gamma_{\mathrm{att}}\big)
        \norm{v_1}_{G_{\mathrm{tr}}}\norm{w_2}_{G_{\mathrm{att}}},
\]
which is negative-definite iff the associated $2\times2$ matrix is positive
definite; minimising the resulting convex condition over $\alpha_1/\alpha_2$ gives
the optimiser \eqref{eq:se3-G} and reduces the condition to \eqref{eq:smallgain}.
Contraction plus the existence of the equilibrium $(e_x,e_v,R_e,\Ome)=(0,0,I,0)$
in the region gives exponential convergence by \cite[Prop.~2.5(i)]{SB13}.
\end{proof}

\begin{remark}[What \ref{H3} costs, and how to discharge it]
\label{rem:gains}
\eqref{eq:smallgain} is a genuine hypothesis, not a corollary. Both gains admit
bounds that make it checkable:
$\gamma_{\mathrm{tr}}$ is governed by
$\partial\Delta/\partial b_3 = -\big(b_3F_{\mathrm{des}}^{\top} + (F_{\mathrm{des}}\cdot b_3)I\big)$,
so $\gamma_{\mathrm{tr}} = O(\sup\norm{F_{\mathrm{des}}})$, i.e.\ large when the
commanded acceleration is large; and $\gamma_{\mathrm{att}}$ is governed by
$\partial b_{3c}/\partial F_{\mathrm{des}} = (I - b_{3c}b_{3c}^{\top})/\norm{F_{\mathrm{des}}}$
together with $\max(k_x,k_v)$, so $\gamma_{\mathrm{att}} = O(\max(k_x,k_v)/F_{\min})$.
Since $\lambda_{\mathrm{att}}$ grows with the attitude gains while
$\gamma_{\mathrm{att}}$ does not, \eqref{eq:smallgain} is always achievable by
making the attitude loop fast --- the usual timescale separation, here with an
explicit certificate rather than an appeal to singular perturbation. Evaluating
the suprema in \eqref{eq:gains} exactly is a finite-dimensional optimisation over
the compact regions; we do not carry it out here, and Theorem~\ref{thm:se3} should
be read as conditional on it.
\end{remark}

%%%%%%%%%%%%%%%%%%%%%%%%%%%%%%%%%%%%%%%%%%%%%%%%%%%%%%%%%%%%%%%%%%%%%%%%%%%%%%%
\section{Region conditions and limitations}
\label{sec:se3-caveats}

\begin{itemize}[leftmargin=1.6em,itemsep=3pt]
\item \textbf{Thrust positivity.} \ref{H2} is not cosmetic. $f\le0$ makes
      $b_{3c}$ undefined or reverses the loop; $\theta_{\max}<\pi/2$ is the
      familiar ``do not flip'' condition and bounds the initial attitude error.
\item \textbf{Rotor saturation.} The mixer inversion produces $f_i$ that must lie
      in $[0,f_{\max}]$. Contraction is a property of the unsaturated vector
      field; on the saturated set the closed loop is a different system and the
      certificate is void. In practice $\mathcal{W}_{\mathrm{tr}}\times\mathcal{W}_{\mathrm{att}}$
      must be intersected with the unsaturated set, and that set must be shown
      forward invariant --- which we have not done.
\item \textbf{Topology.} $T^{*}\SE$ retracts onto $\SE\simeq\Rm P^3\times\Rm^3$,
      not contractible, so as in Part~A no global contraction region exists. The
      attitude error function contributes the same three antipodal critical points.
\item \textbf{Curvature.} The translational block is flat and the coupling is
      handled by \eqref{eq:smallgain}, so curvature enters only through the
      attitude block, where Lemma~\ref{lem:curv}(iii) applies: an inertially
      asymmetric airframe has negative sectional curvatures and hence a
      velocity-bounded $\mathcal{W}_{\mathrm{att}}$.
\item \textbf{What under-actuation actually cost us.} Part~A yields an
      unconditional theorem; Part~B yields a theorem conditional on \ref{H3}. The
      gap is exactly Prop.~\ref{prop:underact}. A direct contraction certificate
      for the under-actuated system on $T^{*}\SE$ --- not decomposed --- would
      require the machinery of kinematic reduction and decoupling vector fields
      \cite[Ch.~8]{BL04}, and as noted in Note~I, \S8, we do not have one.
\end{itemize}

%%%%%%%%%%%%%%%%%%%%%%%%%%%%%%%%%%%%%%%%%%%%%%%%%%%%%%%%%%%%%%%%%%%%%%%%%%%%%%%
\part*{Audit of claims}
\addcontentsline{toc}{section}{Audit of claims}
\setcounter{section}{0}
\renewcommand{\thesection}{C}
\section{}
\label{sec:audit}
\vspace{-2\baselineskip}

\begin{itemize}[leftmargin=1.6em,itemsep=3pt]
\item \emph{Verified symbolically and numerically for this note} (random-sample
      checks to $10^{-10}$, plus closed-loop integration):
      Lemma~\ref{lem:LC} (metric compatibility and torsion-freeness of
      \eqref{eq:LC}); Lemma~\ref{lem:so3-conn} including
      $\Gamma(x,x)=\Jm^{-1}(x\times\Jm x)$; the corrected vee identity
      \eqref{eq:vee-id} --- note the transpose, without which the identity is
      false for non-symmetric $A$; Lemma~\ref{lem:hess} and the Hessian matrix
      \eqref{eq:Hmat}; Prop.~\ref{prop:isotropise}; Lemma~\ref{lem:curv}(i) and
      the $\Gm$-symmetry of $\Jac_v$; the negative sectional curvatures quoted for
      $\Jm=\mathrm{diag}(1,1,8)$; Prop.~\ref{prop:so3-cl} (closed-loop residual
      $2.5\times10^{-16}$ over a $4\,$s time-varying-reference simulation);
      Lemma~\ref{lem:err-kin}; the $\ad^{*}$ formula and Euler--Poincar\'e
      equations \eqref{eq:se3-EP}; Lemma~\ref{lem:tr-err}; and the numbers in
      Remark~\ref{rem:numeric}.
\item \emph{Quoted from Note~I:} Prop.~3.1 (symplectic drift), Prop.~3.2 (no
      symplectic field is contracting), \S4.1 (the $\nabla$-splitting), Prop.~4.3
      (covariant variational equations), Prop.~5.1 (block-diagonal metrics fail),
      Def.~6.1, Assumptions~6.2--6.3, Thm.~6.4 and Remark~6.5, \S5 (curvature and
      topology), \S8 (limitations).
\item \emph{Quoted from Note~II:} Lemma~1.1 (skew-symmetry is $\nabla\Gm=0$),
      Def.~2.1 (transport maps and error functions), Thm.~3.4 (partial contraction).
\item \emph{Quoted from \cite{SB13}:} Def.~2.1--2.2, Thm.~2.3, Prop.~2.5,
      Lem.~3.1 (cascades), Lem.~3.2 and its eq.~(11)--(12) (feedback small gain).
\item \emph{Quoted from \cite{BL04}:} \S4.6.2--4.6.3 (simple mechanical control
      systems), \S5.3 (left-invariant metrics), \S11.1 (PD on Lie groups), Ch.~8
      (kinematic reduction).
\item \emph{Standard, stated without proof:} Koszul formula; Lie--Poisson and
      semidirect-product reduction; Milnor's results on left-invariant metrics
      (Lemma~\ref{lem:curv}(iii) is stated for general $\Jm$ but verified here only
      by example); the structure of the quadrotor mixer.
\item \emph{Assumed, not proved:} \ref{H3} of Theorem~\ref{thm:se3} --- the
      small-gain condition; forward invariance of the unsaturated rotor set;
      $K$-reachability of $\mathcal{W}$ in Theorem~\ref{thm:so3} and of
      $\mathcal{W}_{\mathrm{tr}},\mathcal{W}_{\mathrm{att}}$ in
      Theorem~\ref{thm:se3}. The order-of-magnitude gain estimates in
      Remark~\ref{rem:gains} are heuristic, not bounds.
\item \emph{Attribution:} the controller \eqref{eq:Fdes} and the commanded-attitude
      construction are the standard geometric quadrotor design associated with
      Lee, Leok and McClamroch; the recasting of its interconnection as
      \cite[Lem.~3.2]{SB13} is done here and I have not checked it against the
      literature.
\end{itemize}

%%%%%%%%%%%%%%%%%%%%%%%%%%%%%%%%%%%%%%%%%%%%%%%%%%%%%%%%%%%%%%%%%%%%%%%%%%%%%%%
\begin{thebibliography}{9}

\bibitem{BL04}
F.~Bullo and A.~D.~Lewis,
\emph{Geometric Control of Mechanical Systems}, Texts in Applied Mathematics,
vol.~49, Springer, 2004.

\bibitem{SB13}
J.~W.~Simpson-Porco and F.~Bullo,
\emph{Contraction Theory on Riemannian Manifolds}. Preprint submitted to
\emph{Systems \& Control Letters}, June~5, 2013. Result numbers refer to that
preprint.

\bibitem{NoteI}
\emph{Contraction-Based Control of Simple Mechanical Systems: A Cotangent-Bundle
Formulation} (``Note~I'').

\bibitem{NoteII}
\emph{Partial Contraction for Mechanical Systems, Intrinsically} (``Note~II'').

\end{thebibliography}

\end{document}